\documentclass[12pt,a4paper]{report}

% Packages — compile with pdfLaTeX
\usepackage[T1]{fontenc}
\usepackage[utf8]{inputenc}
\usepackage[lining]{ebgaramond}
\usepackage[cmintegrals,cmbraces]{newtxmath}
\usepackage{microtype}
\usepackage{geometry}
\usepackage{tabularx}
\usepackage{graphicx}
\usepackage{longtable}
\usepackage{booktabs}
\usepackage{array}
\usepackage{listings}
\usepackage{xcolor}
\usepackage{hyperref}
\usepackage{amsmath}
\usepackage{amssymb}
\usepackage{pdflscape}
\usepackage{float}
\usepackage{enumitem}
\usepackage{makecell}
\usepackage{siunitx}
\graphicspath{{./}}

% Page layout
\geometry{
    left=1.8cm,
    right=1.8cm,
    top=2.2cm,
    bottom=2.2cm
}

% Color definitions
\definecolor{crimson}{RGB}{153, 0, 0}
\definecolor{darkblue}{RGB}{0, 0, 153}

% Hyperlink settings
\hypersetup{
    colorlinks=true,
    linkcolor=crimson,
    urlcolor=crimson,
    citecolor=crimson
}

% Allow line breaks at slashes to prevent overfull hbox
\tolerance=2000
\emergencystretch=2em

% Chapter/section title formatting
\usepackage{titlesec}
\titleformat{\chapter}{\normalfont\huge\bfseries}{}{0pt}{}[\vspace{0.3em}]
\titlespacing*{\chapter}{0pt}{0pt}{14pt}
\titleformat{\section}{\normalfont\Large\bfseries}{\thesection}{1em}{}
\titleformat{\subsection}{\normalfont\large\bfseries}{\thesubsection}{1em}{}

\setlength{\parindent}{0pt}

\begin{document}

\pagestyle{plain}

% ─────────────────────────────────────────────
%  Cover Page
% ─────────────────────────────────────────────
\begin{titlepage}
\centering
\vspace*{3cm}

{\fontsize{28}{34}\selectfont\bfseries Deep Learning\par}
\vspace{0.6em}
{\fontsize{20}{26}\selectfont\bfseries Lab 5 \par}
\vspace{0.4em}
{\fontsize{16}{22}\selectfont\bfseries Value-Based Reinforcement Learning \par}

\vspace{1.5cm}
\textcolor{crimson}{\rule{0.6\linewidth}{1.2pt}}
\vspace{1.5cm}

{\large\bfseries Student ID}\par
\vspace{0.4em}
{\large M11407439}\par

\vspace{0.8em}

{\large\bfseries Name}\par
\vspace{0.4em}
{\large MING-HONG, LIN}\par

\vspace{0.8cm}

{\large\bfseries School}\par
\vspace{0.4em}
{\large National Yang Ming Chiao Tung University}\par

\vfill

\textcolor{crimson}{\rule{0.6\linewidth}{0.6pt}}\par
\vspace{0.5em}
{\small Due Date: 2026/05/05}

\end{titlepage}

\tableofcontents
\newpage

\chapter{Introduction}

\chapter{Implementation}
\section{Task 1: DQN on CartPole-v1}

\subsection{Algorithm Overview}

Task~1 implements a \textbf{Deep Q-Network (DQN)} on the CartPole-v1 environment.
The core idea is to approximate the action-value function $Q(s, a;\,\theta)$ with a deep neural
network and learn the optimal policy by minimizing the Bellman error.
DQN introduces three key mechanisms to stabilize training when combining Q-learning with
neural networks:

\begin{itemize}[leftmargin=1.8em]
  \item \textbf{Experience Replay}: At every time step, the transition
        $(s,\,a,\,r,\,s',\,\text{done})$ is stored in a fixed-capacity replay buffer
        (implemented as a \texttt{deque}).
        During training, a mini-batch is sampled uniformly at random, breaking the temporal
        correlation between consecutive samples and improving sample efficiency.

  \item \textbf{Target Network}: A separate network $Q_{\text{target}}(s,a;\,\theta^-)$ with
        frozen parameters is maintained solely for computing Bellman targets.
        Its weights are synchronized from the online network only every $N$ training steps
        (hard copy), preventing the training target from shifting continuously with each
        gradient update and thereby improving convergence stability.

  \item \textbf{$\varepsilon$-greedy Exploration}: With probability $\varepsilon$ the agent
        selects a random action (exploration); with probability $1-\varepsilon$ it selects the
        action with the highest Q value (exploitation).
        $\varepsilon$ decays exponentially during training so the policy gradually shifts from
        exploration toward exploitation.
\end{itemize}

The training loss is the mean squared TD error:

\begin{equation}
  \mathcal{L}(\theta) \;=\;
  \mathbb{E}_{(s,a,r,s')\sim\mathcal{D}}\!\left[
    \Bigl(r + \gamma\,\max_{a'} Q_{\text{target}}(s',a';\,\theta^-) - Q(s,a;\,\theta)\Bigr)^2
  \right]
\end{equation}

\subsection{Implementation Details}

\subsubsection*{Network Architecture}

The state space of CartPole-v1 is a 4-dimensional continuous vector
(cart position, cart velocity, pole angle, pole angular velocity).
A fully connected \textbf{MLP} is used as the Q-network, with the output layer
producing Q-values for each of the two possible actions:

\[
  \text{Linear}(4\!\to\!128) \xrightarrow{\text{ReLU}}
  \text{Linear}(128\!\to\!128) \xrightarrow{\text{ReLU}}
  \text{Linear}(128\!\to\!2)
\]

All weights are initialized with Kaiming Uniform (\texttt{nonlinearity=`relu'}),
and biases are set to zero.

\subsubsection*{Replay Buffer}

A \texttt{deque(maxlen=100{,}000)} is used as a uniform replay buffer.
One transition $(s, a, r, s', \text{done})$ is appended per environment step.
At each training step, a mini-batch of size 32 is drawn via \texttt{random.sample}
to eliminate temporal correlations.

\subsubsection*{$\varepsilon$-greedy Decay}

$\varepsilon$ starts at $1.0$ (pure exploration) and is multiplied by a fixed decay factor
after every call to \texttt{train()}:

\[
  \varepsilon \;\leftarrow\; \varepsilon \times 0.9999925
\]

Decay stops once $\varepsilon$ reaches $\varepsilon_{\min} = 0.01$,
retaining a small amount of exploration throughout training.

\subsubsection*{Loss Computation and Gradient Update}

\begin{enumerate}[leftmargin=2em]
  \item Compute the current Q-value from the online network:
        $Q(s,a;\,\theta)$ via \texttt{gather} on the action dimension.
  \item Compute the TD target with the target network (inside \texttt{torch.no\_grad()}):
        $y = r + \gamma \cdot \max_{a'} Q_{\text{target}}(s',a';\,\theta^-) \cdot (1 - \text{done})$,
        where $\gamma = 0.99$.
  \item Compute the loss:
        $\mathcal{L} = \texttt{nn.functional.mse\_loss}(Q(s,a;\,\theta),\; y)$.
  \item Clip gradients (\texttt{clip\_grad\_norm\_} $\leq 10.0$) and update parameters
        with the Adam optimizer ($\text{lr} = 2.5 \times 10^{-4}$).
\end{enumerate}

\subsubsection*{Target Network Synchronization}

A \textbf{hard update} strategy is adopted: every
\texttt{target\_update\_frequency} $= 1{,}000$ training steps,
all parameters of the online network are copied directly to the target network via
\texttt{load\_state\_dict}.
The target network parameters remain frozen between synchronizations.

\subsubsection*{Replay Start Size}

Training does not begin until the replay buffer contains at least $50{,}000$ transitions.
This warm-up phase ensures sufficient sample diversity in the early mini-batches
and prevents the network from overfitting to a narrow initial distribution of experiences.

\section{Task 2}
\section{Task 3}
\section{Question Answering}

\chapter{Analysis and Discussions}

\chapter{Evaluation Results}

\chapter{Bonus}


\end{document}
